\documentclass[letterpaper,10pt,conference]{ieeeconf}

\IEEEoverridecommandlockouts
\overrideIEEEmargins

\usepackage[T1]{fontenc}
\usepackage{amsmath,amssymb}
\usepackage{graphicx}
\usepackage{booktabs}
\usepackage{array}
\usepackage{multirow}
\usepackage{cite}
\usepackage{url}
\usepackage{flushend}
\usepackage{placeins}
\usepackage[hidelinks]{hyperref}

\newcommand{\etal}{et al.}
\newcommand{\method}{WarpVox}

\title{WarpVox: Observation-Guided Deformation for TSDF Map Correction after Loop Closure}

\author{Anonymous Authors}

\begin{document}
\maketitle
\thispagestyle{empty}
\pagestyle{empty}

\begin{abstract}
Loop closure retroactively changes sensor poses after their range observations have been fused, leaving a TSDF inconsistent with the corrected trajectory. Full reintegration restores consistency but revisits the entire frame history. WarpVox instead caches the range images of SLAM keyframes and rebuilds their associations to the current surface at each trajectory revision. It re-expresses the matched sensor-frame samples with corrected poses to form surface targets, propagates the targets through a single event-level proxy deformation, and reconstructs a Voxblox-compatible TSDF whose zero set follows the warped surface while observed support is primarily inherited from the source volume. The correction stage does not re-integrate measurements: it uses cached keyframe measurements for projective surface re-association, while deformation and reconstruction scale with active geometry. Across 39 Kimera-Multi and 12 KITTI checkpoints at 0.25 m, WarpVox remains within 0.05 and 0.32 percentage points of reintegration at F@25 while reducing mean core time from 109.25 to 2.51 s and from 15.73 to 4.87 s. In a ten-bag final-state surface comparison, a single mesh-only WarpVox update improves mean F@50 by 3.03 percentage points (median 0.20) over Kimera-PGMO, while using $9.27\times$ lower core time than the incremental PGMO mesh deformation required to produce its final mesh. In a single-sequence case study, a geometry-only SplaTAM extension is demonstrated. The supplementary material includes all code, thresholds, and configuration details.
\end{abstract}

\begin{figure*}[t]
\centering
\includegraphics[width=0.92\textwidth]{Figures/teaser_v1.jpg}
\caption{\textbf{WarpVox corrects a TSDF after loop closure without
re-fusing the historical range stream.} The pre-correction map remains locally
coherent but is globally inconsistent with the revised trajectory. WarpVox
re-associates cached keyframe measurements with the current surface, propagates
the resulting targets through one event-level deformation, and reconstructs a
corrected TSDF.}
\label{fig:teaser}
\end{figure*}

% PaperCept keywords: Dense mapping; loop closure; TSDF;
% non-rigid deformation; SLAM.

\section{Introduction}
\label{sec:intro}

Dense 3D maps support navigation, manipulation, and scene understanding in robotics and augmented reality. TSDF methods~\cite{10.1145/3596711.3596726,newcombe2011kinectfusion} integrate range observations incrementally, while sparse volume structures~\cite{10.1145/2508363.2508374,Oleynikova_2017,vizzo2022vdbfusion} make large-scale onboard mapping practical. These systems accept poses from an external SLAM front-end. Loop closure and pose-graph optimization~\cite{Campos_2021,5979949} retroactively correct the trajectory, leaving the TSDF in the old world geometry and therefore inconsistent with the revised poses (Fig.~\ref{fig:teaser}).

% \vspace{1mm}\noindent\textit{The map-rebuild bottleneck.}
The standard remedy is to discard the map and reintegrate all $N$ range frames with corrected poses~\cite{Cramariuc_2023,rosinol2021kimeraslamspatialperception}. Its $O(N|\Omega|)$ cost, for $|\Omega|$ measurements per frame, is paid again after every global correction and can reach minutes on long sessions. A de-integration/reintegration cycle~\cite{han2018flashfusion} avoids a full rebuild but integrates twice for each affected frame; after a global correction nearly every frame is affected, so its cost upper-bounds full reintegration rather than reducing it. We therefore target a narrow interface: consume an already fused TSDF and a revised trajectory, and return a TSDF in the same format that the mapper keeps fusing into.
 

Submap methods such as Voxgraph~\cite{reijgwart2020voxgraph} provide this interface but require the map to be partitioned during fusion, and their rigid per-submap transforms leave intra-submap drift unresolved. The remaining approaches change the representation instead: non-rigid surfel SLAM~\cite{whelan2016elasticfusion} corrects its own tracking state rather than an external TSDF, and neural implicit maps~\cite{sucar2021imap,zhu2022niceslam,liso2024loopyslam,bruns2024neuralgraphmap} require representation-specific re-optimization. Each is a sound design; none returns a corrected TSDF to a mapper already running one.

Our starting point is an invariance: a calibrated range sample expressed in its sensor frame is unchanged by a later world-frame pose correction, although its relation to the current map surface is not. WarpVox therefore uses keyframe range measurements to rebuild projective measurement--surface correspondences at each correction event. Re-expressing the matched samples with corrected poses yields corrected world-space surface targets by projection alone: the correction cost then follows the current map rather than the number of measurements behind it or the volumetric work of fusing them.

WarpVox combines four stages: (i) projectively re-associate keyframe range measurements with the current pre-correction surface, (ii)~compute robust targets from the corrected poses, (iii)~propagate sparse displacements through a single ARAP solve~\cite{ARAP_modeling:2007} on a QEM proxy, and (iv)~reconstruct a TSDF whose zero set follows the warped surface while observed support, sign, and weight primarily follow the source volume. At each loop-closure, this corrected volume replaces the inconsistent active TSDF; Voxblox then resumes fusion from the updated state until a later trajectory revision triggers the same correction procedure again. Only Stages~1--2 use keyframe measurements to construct correction targets; deformation and reconstruction operate on the active geometry and output volume.

\begin{itemize}
% \item We introduce an event-level system that converts an external trajectory revision and an already fused active TSDF into a corrected, fusion-ready TSDF without replaying the complete historical range stream.
% \item The system constructs correction targets through projective measurement--surface reassociation, propagates them through one proxy deformation, and reconstructs source-supported TSDF fields so that the corrected volume can replace the active map.
\item An observation-derived target construction that turns cached keyframe range measurements and a revised trajectory into corrected surface targets by projection, without replaying volumetric integration. 
\item A source-supported TSDF reconstruction that pairs the warped zero set with support, sign, and integration weights pulled back from the source volume, returning a native Voxblox map that replaces the active TSDF and continues fusing.
\item An evaluation of RGB-D and LiDAR correction against reintegration and mapping baselines over 51 loop-closure cases, together with a separate comparison of final-state surface correction against Kimera-PGMO on ten complete bags.
\end{itemize}

\input{sections/introduction_addition}
\section{Related Work}
\label{sec:related}

\vspace{1mm}\noindent\textit{Volumetric mapping.}
TSDF fusion~\cite{10.1145/3596711.3596726,newcombe2011kinectfusion} and sparse volume structures~\cite{10.1145/2508363.2508374,Oleynikova_2017,vizzo2022vdbfusion} enable incremental dense mapping from externally supplied poses. These systems normally assume that the input trajectory remains fixed and do not expose a mechanism for revising a persistent volume after a pose-graph backend changes historical poses.

\vspace{1mm}\noindent\textit{Dense-map correction after loop closure.}
Kintinuous~\cite{whelan2015kintinuous} deforms accumulated point-cloud slices
rather than returning a corrected persistent TSDF, and
ElasticFusion~\cite{whelan2016elasticfusion} repeatedly deforms a surfel map
inside its own tracking system. Voxgraph~\cite{reijgwart2020voxgraph} instead
optimizes rigid submaps, leaving spatially varying intra-submap errors
unresolved. Kimera-PGMO~\cite{kimera_pgmo} jointly optimizes robot poses and a simplified timestamped mesh in a unified deformation graph, then displaces each vertex using its $k$ nearest optimized control nodes within a time tolerance. In Kintinuous and Kimera-PGMO the displacement of a map element is therefore set by its spatio-temporal position in the graph rather than by the measurement that produced it. We report the latter separately (Sec.~\ref{sec:pgmo}) because it produces its own trajectory rather than consuming a fixed corrected one.

\vspace{1mm}\noindent\textit{Surface deformation.}
ARAP~\cite{ARAP_modeling:2007} and embedded deformation~\cite{10.1145/1276377.1276478} preserve local shape while satisfying sparse controls. DynamicFusion~\cite{7298631}, KillingFusion~\cite{slavcheva2017killingfusion}, and VolumeDeform~\cite{innmann2016volumedeform} estimate continuous non-rigid scene motion. WarpVox instead addresses a static scene with a discrete trajectory revision: corrected poses generate the controls, and a single QEM-proxy deformation propagates them over the extracted surface.

\vspace{1mm}\noindent\textit{Neural and Gaussian maps.}
Neural SLAM systems~\cite{sucar2021imap,zhu2022niceslam,cui2023ngelslam} optimize implicit scene representations; Loopy-SLAM~\cite{liso2024loopyslam} revisits that optimization after loop closure, while Neural Graph Map~\cite{bruns2024neuralgraphmap} anchors local fields to pose-graph nodes. SplaTAM~\cite{keetha2024splatam} maps with 3D Gaussians~\cite{kerbl2023gaussians}, and LoopSplat~\cite{zhu2025loopsplat} aligns Gaussian submaps. These representations do not generally provide the explicit TSDF/ESDF observed-space semantics required by robotic distance-map pipelines. WarpVox therefore focuses on TSDF/ESDF pipelines; its Gaussian experiment isolates only the transferability of observation-derived targets.

\input{sections/related_work_addition}
\section{Method}
\label{sec:method}

\noindent\textbf{Assumptions.}
Following the local-consistency assumption of
Voxgraph~\cite{reijgwart2020voxgraph}, we assume that TSDF fusion under the
drifted pre-correction trajectory remains locally valid: surface geometry and
connectivity are coherent in local neighborhoods although the map is globally
inconsistent. WarpVox corrects this global inconsistency after loop closure;
severely noisy or topologically corrupted source maps are outside its scope.

\subsection{Problem and Correction Inputs} Let $T^-(t),T^+(t)\in SE(3)$ map the sensor frame at time $t$ to the world frame before and after trajectory correction. The map built with $T^-$ consists of a TSDF $\Phi^-$, its weight field $w^-$, voxel size $s$, and truncation distance $\tau$. During online fusion, Voxblox continuously maintains the Marching-Cubes surface $\mathcal M^-=(V^-,F)$ of the active TSDF~\cite{10.1145/37402.37422}. On a trajectory correction, this mesh becomes the source surface. Let $\mathcal A$ denote the active TSDF cells that contribute triangles to $\mathcal M^-$. WarpVox stores the timestamped, modality-native range images of SLAM keyframes,
\begin{equation}
\mathcal C_{\rm KF}=\{k_i=(t_i,D_i):i\in{\rm KF}\},
\end{equation}
where ${\rm KF}$ indexes the retained SLAM keyframes and $D_i$ is an RGB-D depth image or a spherical LiDAR range image. Together with $T^-$, these measurements form the observation input to Stage~1. The measurements remain available across correction events. Fig.~\ref{fig:pipeline} summarizes the four stages.

\begin{figure}[t]
  \centering
  \includegraphics[width=\linewidth]{Figures/pipeline_v5.jpg}
  \caption{WarpVox extracts the current source surface from the fused TSDF, re-establishes projective correspondences with keyframe measurements, propagates the resulting targets through one proxy deformation per correction, and reconstructs a source-supported TSDF.}
  \label{fig:pipeline}
\end{figure}

\subsection{Stage 1: Keyframe--Surface Association}

Stage~1 projects the current surface into the keyframe range images and returns vertex-indexed observation records,
\begin{equation}
\begin{aligned}
\mathcal O&=\mathcal R(\mathcal M^-,T^-,\mathcal C_{\rm KF}),\\
\mathcal O_j&=\{(r,\alpha_{rj})\}.
\end{aligned}
\label{eq:resolved_records}
\end{equation}
where each record $r$ identifies a keyframe range image and a valid projected range sample, $t_r$ is its acquisition time, and $\alpha_{rj}\in[0,1]$ is the
effective reliability of associating it with source vertex $v_j^-$. For target construction, each vertex retains at most 15 association records, giving fixed-size per-vertex state. RGB-D restricts candidates with a camera-centre KD-tree, LiDAR with a
sensor-position grid and an 80\,m radius check. The correction
backend therefore receives vertex-indexed records and never re-fuses
measurements into the volume. Its contract is
\begin{equation}
(\Phi^-,w^-,\mathcal M^-,T^+,\mathcal O)
\longmapsto (\Phi^+,w^+),
\label{eq:contract}
\end{equation}
where the warped mesh is intermediate and $(\Phi^+,w^+)$ is a
native Voxblox TSDF volume.

For an RGB-D record $r$, pixel $\bar u_r=[u_r,v_r,1]^\top$ with measured
z-depth $z_r$ defines the invariant sensor-frame sample
\begin{equation}
m_r=z_rK^{-1}\bar u_r,
\qquad
\widetilde m_r=[m_r^\top,1]^\top.
\label{eq:sensor_sample}
\end{equation}
The inverse pre-correction pose projects source vertex $v_j^-$ into the record frame. Candidate agreement combines depth residual, three-dimensional position, and viewing direction. Let $n_j^-$ be the source normal, $\ell_{rj}^-$ the unit viewing direction in the pre-correction world frame, and $z_{rj}^{\rm proj}$ the projected vertex depth. We define
\begin{equation}
\begin{aligned}
e^g_{rj}
&=\left\|[T^-(t_r)\widetilde m_r]_{1:3}-v_j^-\right\|,\\
e^z_{rj}
&=z_{rj}^{\rm proj}-z_r,\\
g_{rj}
&=\left|(n_j^-)^\top \ell_{rj}^-\right|.
\end{aligned}
\label{eq:association_residuals}
\end{equation}
The effective association reliability is
\begin{equation}
\alpha_{rj}=
\exp\!\left(-\frac{e^g_{rj}}{0.048\,\mathrm{m}}\right)
\exp\!\left(-\frac{|e^z_{rj}|}{0.08\,\mathrm{m}}\right)
\frac{g_{rj}}{1+z_{rj}^{\rm proj}/8\,\mathrm{m}}.
\label{eq:association_weight}
\end{equation}
RGB-D candidates satisfy the depth-dependent geometric gate $e^g_{rj}\le 0.08\,\mathrm{m}+0.01\,z_{rj}^{\rm proj}$, and candidates with $\alpha_{rj}<0.05$ are rejected. LiDAR uses spherical projection with radial range $\rho$: candidates satisfy $e^g_{rj}\le 0.25\,\mathrm{m}+0.003\,\rho_{rj}^{\rm proj}$ and are scored by \begin{equation} \alpha_{rj}=\exp\!\left(-\frac{e^g_{rj}}{0.15\,\mathrm{m}}\right) \frac{g_{rj}}{1+\rho_{rj}^{\rm proj}/40\,\mathrm{m}}, \label{eq:association_weight_lidar} \end{equation} with no further $\alpha$ threshold. This establishes geometric support, not exact voxel-integration provenance. The attenuation in \eqref{eq:association_weight} ranks candidates during association; the separate factor in \eqref{eq:target_initial} models target reliability after association.

\subsection{Stage 2: Corrected Surface Targets}

After loop closure, the timestamp $t_r$ selects the corresponding corrected pose. Missing interior poses use linear translation interpolation and SLERP rotation over gaps of at most three times the median trajectory spacing~\cite{shoemake1985slerp}. The invariant sample then defines the corrected candidate
\begin{equation}
q_r^+=\left[T^+(t_r)\widetilde m_r\right]_{1:3}.
\label{eq:corrected_target}
\end{equation}
This candidate is geometrically consistent when the static-scene association, calibration, and
corrected pose are valid. Let
$\mathcal J_j=\{r:(r,\alpha_{rj})\in\mathcal O_j\}$. We first compute
\begin{equation}
\omega_{rj}=\frac{\alpha_{rj}}{1+z_r/3\,\mathrm{m}},
\qquad
\bar q_j=
\frac{\sum_{r\in\mathcal J_j}\omega_{rj}q_r^+}
     {\sum_{r\in\mathcal J_j}\omega_{rj}}.
\label{eq:target_initial}
\end{equation}
For LiDAR, $\omega_{rj}=\alpha_{rj}/(1+\|q_r^+-v_j^-\|/20\,\mathrm{m})$ and the floor in \eqref{eq:target_inliers} is $0.05\,\mathrm{m}$.
The reliability-weighted center is paired with an unweighted median residual
scale so that no single high-weight candidate controls both the estimate and
its rejection threshold:
\begin{equation}
\begin{aligned}
\theta_j&=\max\!\left(
0.02\,\mathrm{m},
2.5\operatorname*{median}_{r\in\mathcal J_j}
\|q_r^+-\bar q_j\|\right),\\
I_j&=\left\{r\in\mathcal J_j:
\|q_r^+-\bar q_j\|\le\theta_j\right\}.
\end{aligned}
\label{eq:target_inliers}
\end{equation}
When fewer than three candidates are available, $I_j=\mathcal J_j$. The final target is
\begin{equation}
v_j^*=
\frac{\sum_{r\in I_j}\omega_{rj}q_r^+}
     {\sum_{r\in I_j}\omega_{rj}}.
\label{eq:target}
\end{equation}
Each vertex with at least one valid record yields a target; the resulting controls are screened by the robust spatial-consistency filter of Stage~3.

\subsection{Stage 3: Surface Deformation}

The deformation stage propagates observation-derived motion; it does not
infer correction from shape alone. Let $d_j=v_j^*-v_j^-$ be a control
displacement. Before deformation, control-aware mesh cleaning preserves
connected components carrying observation-derived controls. A fixed
component-aware robust filter then removes isolated controls whose
displacements disagree with the surrounding accepted motion. Small
disconnected components retained only because they initially carried controls
are excluded if all such controls are subsequently rejected. Remaining
components with no surviving control inherit a robust local translation from
nearby accepted controls through synthetic constraints; these constraints
define otherwise underconstrained motion but are not treated as observations.

Component-separated spatial clustering~\cite{rossignac1993vertexclustering}
constructs a triangular proxy $\mathcal P^-=(P^-,F_P)$ with a target budget
of approximately 60K vertices and selects one QEM
representative~\cite{garland1997qem} per occupied cell. Let
$\pi:V^-\rightarrow P^-$ map source vertices to representatives and
$\mathcal K_a=\{j:\pi(j)=a\}$ collect their controls. For
$\mathcal K_a\neq\emptyset$, the initial proxy target is
\begin{equation}
p_a^{(0)}=p_a^-+
\frac{1}{|\mathcal K_a|}
\sum_{j\in\mathcal K_a}d_j.
\label{eq:proxy_target}
\end{equation}
A robust local-rigid consensus
~\cite{schoenemann1966procrustes,yao2020welsch} regularizes these proxy
targets. An observation-derived target is replaced by the fitted motion only
when a strict majority of the evaluated targets in the same connected
component are inconsistent with that fit; isolated disagreements remain
unchanged. Synthetic targets are regularized by the same local motion field.
Let $\mathcal C_P\subseteq\{1,\ldots,|P^-|\}$ index the resulting measured
or synthetic proxy targets $p_a^*$.

WarpVox performs one proxy ARAP solve per correction~\cite{ARAP_modeling:2007, 10.1007/978-3-662-05105-4_2,libigl}:
\begin{equation}
\begin{aligned}
\min_{P^+,\{R_a\in SO(3)\}}\quad
&\frac{1}{2}\sum_{a=1}^{|P^-|}
\sum_{b\in\mathcal N_P(a)}c_{ab}\\[-1mm]
&\left\|(p_a^+-p_b^+)-R_a(p_a^--p_b^-)\right\|^2,\\
\text{s.t.}\quad
&p_a^+=p_a^*,\qquad a\in\mathcal C_P .
\end{aligned}
\label{eq:arap}
\end{equation}
Here $\mathcal N_P(a)$ is the one-ring neighbourhood induced by $F_P$, and $c_{ab}=c_{ba}$ are symmetric cotangent-derived weights; the factor $1/2$ removes double counting. Once accepted, controls are imposed exactly so that ARAP does not attenuate the supplied trajectory correction. The standard local--global iterations solve this single optimization.

Proxy displacement is transferred to $V^-$ by topology-local barycentric
interpolation. Let $f_j$ be the closest triangle to $v_j^-$ in
$\operatorname{star}(\pi(j))$. With barycentric coordinates $\lambda_{jb}$ of
its closest point,
\begin{equation}
v_j^+=v_j^-+\sum_{b\in f_j}\lambda_{jb}(p_b^+-p_b^-),
\qquad \lambda_{jb}\ge0,\quad\sum_{b\in f_j}\lambda_{jb}=1.
\label{eq:proxy_transfer}
\end{equation}
The ARAP solve and topology-local transfer modify vertex positions only;
before the safety pass, they preserve the source connectivity and the
separation of disconnected components. A subsequent safety pass harmonically
relaxes regions exhibiting severe deformation-induced distortion. Faces whose
edge or area expansion remains unresolved are removed, but no new connections
are created and the surface is not retriangulated. The retained source and
warped meshes therefore share the same indexed face subset,
\begin{equation}
\bar{\mathcal M}^-=(\bar V^-,\bar F),\qquad
\mathcal M^+=(\bar V^+,\bar F),\qquad
\bar F\subseteq F.
\label{eq:retained_topology}
\end{equation}

\subsection{Stage 4: Source-Supported TSDF Reconstruction}

The warped surface supplies the corrected zero set
~\cite{sussman1999redistancing}, while the source TSDF supplies observed
support, sign orientation, integration weights, and free space. Let $n_b$ be
the number of voxels per block edge, $h=n_bs$ the block width, $o_b=hb$ the
origin of block $b$, and
\begin{equation}
\operatorname{vox}(b)=
\left\{o_b+s\left(k+\tfrac12\mathbf 1\right):
k\in\{0,\ldots,n_b-1\}^3\right\}
\label{eq:block_voxels}
\end{equation}
its voxel centres. Let $\mathcal B_{\rm fs}^+$ denote the additional blocks
reached by forward transport of observed source free space. With
$\mathcal N_1$ denoting a one-block Chebyshev dilation, the final allocation
domain is
\begin{equation}
\begin{aligned}
\mathcal B_V^+&=\{\lfloor \bar v/h\rfloor:\bar v\in\bar V^+\},\\
\mathcal B^+&=\mathcal N_1(\mathcal B_V^+)\cup\mathcal B_{\rm fs}^+,\\
\mathcal X^+&=\bigcup_{b\in\mathcal B^+}\operatorname{vox}(b).
\end{aligned}
\label{eq:allocation}
\end{equation}
No scene-wide bounding box is voxelized: the output contains the
warped-surface footprint and only the additional blocks reached by transported
free-space support, giving
$|\mathcal X^+|\le
n_b^3\!\left(27|\mathcal B_V^+|+|\mathcal B_{\rm fs}^+|\right)$.
For $x\in\mathcal X^+$, a closest-triangle
query~\cite{10.1145/2601097.2601199} on
$\mathcal M^+$ returns face $f$, closest point $q^+$, and barycentric
coordinates $\beta$. Shared face indices locate the corresponding source point
\begin{equation}
q^-=\sum_{k\in f}\beta_k\bar v_k^-.
\label{eq:source_correspondence}
\end{equation}
Let
\begin{equation}
R_f^\pm=
\begin{bmatrix}
t_{f,1}^\pm & t_{f,2}^\pm & n_f^\pm
\end{bmatrix}\in SO(3)
\label{eq:triangle_frames}
\end{equation}
be consistently oriented orthonormal frames of the corresponding source and
warped triangles. The first tangent follows the same indexed edge and
$t_{f,2}^\pm=n_f^\pm\times t_{f,1}^\pm$. A piecewise local-rigid pullback,
motivated by deformation transfer~\cite{sumner2004deformationtransfer}, maps
$x$ to the source volume:
\begin{equation}
\Psi_f(x)=q^-+R_f^-(R_f^+)^\top(x-q^+),
\qquad x^-=\Psi_f(x).
\label{eq:inverse_map}
\end{equation}
This preserves tangential and normal offsets relative to the corresponding
triangles.
With source TSDF samples $\widetilde\Phi^-$ and $\widetilde w^-$, let
\begin{equation}
r(x)=\|x-q^+\|,
\qquad
d^+(x)=\min\!\left(\tau,r(x)\right).
\label{eq:warped_distance}
\end{equation}
A complete trilinear source sample is valid only when every source voxel with
a nonzero interpolation coefficient carries positive weight. When such a
sample is available, WarpVox inherits its interpolated weight and its sign
evidence. When a complete distance sample is unavailable but the pullback
still reaches observed source support, only the source weight is inherited;
the corrected distance remains defined by the warped surface. Let
$w_{\rm src}^+(x)$ denote the resulting inherited source weight.

Let $s_f^+(x)$ denote the oriented warped-surface sign. Whenever possible,
the corresponding source TSDF orients the warped normal using its variation
along the source-triangle normal. If this orientation is unavailable, a
sign-consistent source sample is used in the surface-local band; otherwise
the oriented warped-surface sign is retained. We denote the resulting sign by
$s_{\rm tr}(x)$.

The initial source-supported candidates are
\begin{equation}
\begin{aligned}
\widehat{\Omega}_{\rm tr}
&=\left\{x\in\mathcal X^+:
r(x)\le\tau,\ w_{\rm src}^+(x)>0\right\},\\
\Omega_{\rm n}
&=\left\{x\in\widehat{\Omega}_{\rm tr}:r(x)<1.5s\right\},\\
\widehat{\Omega}_{\rm e}
&=\widehat{\Omega}_{\rm tr}\setminus\Omega_{\rm n}.
\end{aligned}
\label{eq:transferred_support}
\end{equation}

Let $\partial_{\rm open}\mathcal M^+$ denote the union of open boundary edges
and their incident vertices. Where source support is unavailable, unsupported
voxels in the same finite surface band receive a unit fallback prior:
\begin{equation}
\Omega_{\rm fb}=
\left\{x\in\mathcal X^+\setminus\widehat{\Omega}_{\rm tr}:
r(x)<1.5s,\ q^+\notin\partial_{\rm open}\mathcal M^+\right\}.
\label{eq:fallback_support}
\end{equation}
Its weight is
\begin{equation}
w_{\rm fb}(x)=
\operatorname{clamp}_{[0,1]}\!\left(
1.5-\frac{r(x)}{s}\right).
\label{eq:fallback_weight}
\end{equation}
This fallback is a finite prior required to serialize and continue updating
the corrected zero set; it is not interpreted as recovered observation count.

From $\widehat{\Omega}_{\rm e}$, WarpVox retains only candidates that are
sign-consistent with neighbouring active support and connected to the
surface-local support $\Omega_{\rm n}\cup\Omega_{\rm fb}$. Let the retained
set be $\Omega_{\rm e}$ and define
$\Omega_{\rm tr}=\Omega_{\rm n}\cup\Omega_{\rm e}$. For
$x\in\Omega_{\rm tr}$,
\begin{equation}
\Phi_{\rm tr}^+(x)=s_{\rm tr}(x)d^+(x),
\qquad
w_{\rm tr}^+(x)=w_{\rm src}^+(x).
\label{eq:transported_tsdf}
\end{equation}
To preserve observed free space beyond the surface-local bands, confidently
positive source TSDF voxels are grouped locally and transported forward using
the paired source and warped surface frames. For a local group $c$, let
$q_c^-$ and $q_c^+$ be corresponding points on the paired triangles and let
$R_c^-$ and $R_c^+$ be their local frames. The forward map is
\begin{equation}
\Gamma_c(y)=
q_c^+ + R_c^+(R_c^-)^\top(y-q_c^-).
\label{eq:forward_free_space}
\end{equation}
Mapped samples may allocate additional output blocks. They are written only
into previously unsupported voxels and are rejected when they would introduce
a new zero crossing adjacent to retained occupied support. The accepted
samples define $\Omega_{\rm fs}$ and retain positive, source-derived TSDF
values and weights $(\Phi_{\rm fs}^+,w_{\rm fs}^+)$.

Define
$\Omega^+=\Omega_{\rm tr}\cup\Omega_{\rm fb}\cup\Omega_{\rm fs}$.
The branches are disjoint because fallback and free-space transport write only
to unsupported voxels. The reconstructed TSDF is
\begin{equation}
\Phi^+(x)=
\begin{cases}
\Phi_{\rm tr}^+(x), & x\in\Omega_{\rm tr},\\
s_f^+(x)d^+(x), & x\in\Omega_{\rm fb},\\
\Phi_{\rm fs}^+(x), & x\in\Omega_{\rm fs},
\end{cases}
\label{eq:hybrid_tsdf}
\end{equation}
with weights
\begin{equation}
w^+(x)=
\begin{cases}
w_{\rm tr}^+(x), & x\in\Omega_{\rm tr},\\
w_{\rm fb}(x), & x\in\Omega_{\rm fb},\\
w_{\rm fs}^+(x), & x\in\Omega_{\rm fs},\\
0, & \text{otherwise}.
\end{cases}
\label{eq:hybrid_weight}
\end{equation}
Because each branch assigns positive weight on its support, by construction
$\Omega^+=\{x\in\mathcal X^+:w^+(x)>0\}$.

When no source TSDF is available, the final Stage-3 surface is converted into
a closest-triangle signed Euclidean distance field within $1.5$ voxel widths,
using the same finite fallback prior. The resulting $(\Phi^+,w^+)$ is
serialized in the native Voxblox TSDF format; its MC mesh is used for surface
evaluation and visualization. The returned volume replaces the active Voxblox
TSDF, so subsequent fusion and later correction events proceed from the
corrected state.

\subsection{Complexity and Scope} Let $A=|\mathcal A|$, $K_{\rm KF}=|\mathcal C_{\rm KF}|$, $P=|P^-|$, and $X^+=|\mathcal X^+|$ be the numbers of active source cells, stored keyframe range images, proxy vertices, and allocated output voxels. At fixed resolution, MC yields a constant number of surface elements per active cell, so $|V^-|=O(A)$; storage is $O(\sum_{i\in{\rm KF}}|D_i|)$ for the measurements and $O(A)$ for the per-vertex association state. With $\kappa\le K_{\rm KF}$ the maximum number of keyframes returned for one source vertex, association costs $O(K_{\rm KF}\log K_{\rm KF}+A(\log K_{\rm KF}+\kappa))$ for RGB-D and $O(K_{\rm KF}+A(1+\kappa))$ for LiDAR under the standard output-sensitive index model, giving \begin{equation} C_{\rm corr}=C_{\rm assoc}+C_{\rm ARAP}(P,F_P)+C_{\rm vol}(X^+,\bar F), \label{eq:complexity} \end{equation} where $C_{\rm vol}$ is the Stage-4 volumetric reconstruction cost. The block-local allocation in \eqref{eq:allocation} makes $X^+$ follow the warped-surface footprint and transported source support rather than scene extent, and the proxy targets approximately 60K vertices. The retained association state per surface element is bounded by a constant, and each output voxel is written once rather than once per contributing measurement. Full reintegration instead re-fuses every historical measurement at $O(N|D|)$. The two costs are indexed by different quantities: the keyframes and the active map for WarpVox, the measurement history for reintegration.

\section{Experiments}
\label{sec:experiments}

\subsection{Evaluation Protocol}

\begin{table}[t]
\centering
\caption{\textbf{Mean reconstruction accuracy and core time at
0.25 m.}}
\label{tab:broad_main}
\scriptsize
\setlength{\tabcolsep}{1.7pt}
\begin{tabular}{llrrrr}
\toprule
Data & Method & F@10 (\%) & F@25 (\%) & F@50 (\%) & Time (s) $\downarrow$ \\
\midrule
\multirow{6}{*}{Kimera (39)}
 & Uncorrected & 7.00 & 19.29 & 34.91 & -- \\
 & Open3D & 8.13 & 22.94 & 41.30 & -- \\
 & Voxgraph & 8.27 & 22.78 & 39.49 & 5.54 \\
 & Wavemap & 1.73 & 11.54 & 30.91 & 8.32 \\
 & Voxblox reint. & 10.93 & 28.71 & 48.67 & 109.25 \\
 & \textbf{WarpVox} & 11.13 & 28.76 & 48.39 & 2.26 \\
\midrule
\multirow{6}{*}{KITTI (12)}
 & Uncorrected & 7.13 & 27.91 & 49.22 & -- \\
 & Open3D & 13.50 & 38.58 & 61.02 & -- \\
 & Voxgraph & 9.38 & 31.89 & 52.49 & 4.53 \\
 & Wavemap & 5.23 & 33.10 & 62.55 & 2.68 \\
 & Voxblox reint. & 13.21 & 44.29 & 67.45 & 15.73 \\
 & \textbf{WarpVox} & 13.00 & 43.97 & 67.07 & 4.87 \\
\bottomrule
\end{tabular}
\parbox{\columnwidth}{\vspace{2pt}\scriptsize A dash marks an untimed update: the uncorrected map applies no correction, and Open3D was not timed.}
\end{table}

\noindent\textbf{Data and execution.}
Each bag is processed once end-to-end; loop-closure checkpoints are not
replayed independently. ORB-SLAM3 and Voxblox run from the start, and each
WarpVox correction replaces the active volume before fusion continues.
Required keyframe range images are cached online, so loop events require no
bag rewind or separate preprocessing. Associations, targets, and deformation
state are discarded after each correction; the cache and corrected TSDF
persist. We evaluate 39 RGB-D loop closures from 18 Kimera-Multi
sequences~\cite{tian23arxiv_kimeramultiexperiments} and 12 LiDAR cases from
six KITTI sequences~\cite{Geiger2013IJRR}. The uncompressed 16-bit RGB-D
cache occupies 0.9 -- 3.6\,GiB (mean 2.1\,GiB over 1349 -- 5707 keyframes).

\noindent\textbf{Reference and alignment.}
All methods reconstruct the same temporal interval. Kimera-Multi uses the
official ground-truth point cloud cropped to a 5\,m tube around the GT
trajectory, matching Voxblox's maximum integration range. KITTI uses the
Velodyne cloud aggregated at ground-truth poses within the mapper's 25\,m
ray range. Predictions outside the GT frame are aligned using only
timestamp-matched trajectory positions; no mesh-to-GT registration is used.
Voxgraph uses the trajectory recovered from its optimized submap-pose graph,
while reconstructions driven by the corrected ORB-SLAM3 trajectory share one
alignment per case.

\noindent\textbf{Baselines.}
Uncorrected Voxblox uses the pre-correction trajectory.
Open3D~\cite{zhou2018open3dmodernlibrary3d} provides an independent TSDF
rebuild, Wavemap~\cite{reijgwart2023wavemap} a wavelet occupancy rebuild,
and Voxblox reintegration~\cite{Oleynikova_2017} the representation-matched
TSDF reference. These methods re-fuse every range frame using the corrected
trajectory interpolated to each timestamp. Voxgraph~\cite{reijgwart2020voxgraph}
optimizes rigid TSDF submaps in its independently optimized frame. WarpVox
starts from the active uncorrected TSDF and uses keyframe measurements only
for target construction.

\noindent\textbf{Metrics.}
We report F@10/25/50 from bidirectional nearest-neighbour tests at
0.10/0.25/0.50\,m. All methods share each case's crop, two-million-point
sample, random seed, and thresholds. Main TSDF experiments use 0.25\,m
voxels and 1.0\,m truncation; Wavemap uses a matched 0.25\,m finest level
and is evaluated on its official occupancy point cloud. The 11-case
high-resolution subset and ten-bag subset were each sampled once; the latter
is shared by the PGMO comparison and ablation.

\noindent\textbf{Timing.}
We report in-memory core map-update time on an AMD Ryzen 7 9700X with
64\,GB RAM. Depending on the method, this includes TSDF or occupancy
integration, submap or mesh optimization, or the applicable WarpVox stages.
Disk I/O and PGMO MeshFrontend construction are excluded.

\subsection{Reconstruction Fidelity and Update Cost}

At 0.25\,m, WarpVox is the closest method to reintegration at every
threshold (Table~\ref{tab:broad_main}).

On Kimera-Multi, WarpVox differs from reintegration by
$+0.20/+0.05/-0.28$ pp at F@10/F@25/F@50. WarpVox is $48.43\times$ faster
than reintegration on Kimera-Multi and $3.23\times$ on KITTI. The smaller
KITTI gain is consistent with lower-cost sparse-LiDAR replay and WarpVox's
larger active surface and output volume. Voxgraph improves over uncorrected
Voxblox at every threshold by aligning rigid SDF submaps, but stays $5.98$
and $12.08$ pp below WarpVox at F@25 on the two datasets, consistent with
uncorrected intra-submap drift. Wavemap shows low F@10 but competitive KITTI
F@50. This is consistent with a hierarchical occupancy representation whose
0.25\,m setting is its finest level rather than a uniform map resolution;
its grid-centred evaluation points can also sit up to 0.22\,m from the true
surface.

At 0.10\,m on 11 Kimera-Multi cases, WarpVox obtains 17.13/39.77/57.83 at
F@10/25/50 versus 17.21/39.68/57.12 for reintegration, while reducing core
time from 192.548 to 9.207\,s ($20.91\times$).
Table~\ref{tab:runtime_breakdown} separates association, target estimation,
deformation, topology processing, and TSDF reconstruction at both
resolutions. Figure~\ref{fig:qualitative_ral} complements the quantitative
comparison.

\begin{table}[t]
\centering
\caption{\textbf{Mean WarpVox core-time breakdown in seconds.}}
\label{tab:runtime_breakdown}
\scriptsize
\setlength{\tabcolsep}{1.25pt}
\begin{tabular}{lrrrrrr}
\toprule
Setting & Assoc. & Target & ARAP & Topol. & TSDF & Total \\
\midrule
Kimera, 25 cm (39) & 0.199 & 0.399 & 1.419 & 0.008 & 0.231 & 2.256 \\
KITTI, 25 cm (12)   & 0.310 & 0.069 & 2.457 & 0.147 & 1.888 & 4.872 \\
Kimera, 10 cm (11) & 2.623 & 0.244 & 5.298 & 0.004 & 1.038 & 9.207 \\
\bottomrule
\end{tabular}
\vspace{1pt}
\end{table}

\begin{figure*}[!t]
\centering
\setlength{\tabcolsep}{0pt}
\renewcommand{\arraystretch}{0.2}
\newcommand{\qa}[1]{\includegraphics[width=.245\linewidth,trim=150 150 150 150,clip]{#1}}
\newcommand{\qb}[1]{\includegraphics[width=.245\linewidth,trim=250 70 250 70,clip]{#1}}
\newcommand{\qc}[1]{\includegraphics[width=.245\linewidth,trim=200 20 200 20,clip]{#1}}
\newcommand{\qd}[1]{\includegraphics[width=.245\linewidth,trim=100 20 100 20,clip]{#1}}
\begin{tabular}{cccc}
\qa{Figures/qualitative/5_voxblox_pre.png} &
\qa{Figures/qualitative/5_voxgraph_post_sep.png} &
\qa{Figures/qualitative/5_warped_tsdf.png} &
\qa{Figures/qualitative/5_voxblox_post.png} \\[2pt]
\qb{Figures/qualitative/28_voxblox_pre.png} &
\qb{Figures/qualitative/28_voxgraph_post_sep.png} &
\qb{Figures/qualitative/28_warped_tsdf.png} &
\qb{Figures/qualitative/28_voxblox_post.png} \\[2pt]
\qc{Figures/qualitative/40_voxblox_pre.png} &
\qc{Figures/qualitative/40_voxgraph_post_sep.png} &
\qc{Figures/qualitative/40_warped_tsdf.png} &
\qc{Figures/qualitative/40_voxblox_post.png} \\[2pt]
\qc{Figures/qualitative/76_voxblox_pre.png} &
\qc{Figures/qualitative/76_voxgraph_post_sep.png} &
\qc{Figures/qualitative/76_warped_tsdf.png} &
\qc{Figures/qualitative/76_voxblox_post.png} \\[2pt]
\qd{Figures/qualitative/77_voxblox_pre.png} &
\qd{Figures/qualitative/77_voxgraph_post_sep.png} &
\qd{Figures/qualitative/77_warped_tsdf.png} &
\qd{Figures/qualitative/77_voxblox_post.png} \\[2pt]
\small Voxblox$_{\mathrm{Pre}}$ &
\small Voxgraph &
\small \textbf{WarpVox} &
\small Voxblox reintegration
\end{tabular}
\caption{\textbf{Qualitative reconstruction comparison at 0.25 m.} Three
Kimera-Multi loop-closure cases are shown in the top three rows and two KITTI
cases in the bottom two rows; each row uses the same view across methods.
Colors encode vertex height.}
\label{fig:qualitative_ral}
\end{figure*}

\subsection{Comparison with Kimera-PGMO}
\label{sec:pgmo}

This experiment compares Kimera-PGMO's deformation-graph correction with
WarpVox's observation-derived targets and proxy ARAP. Both use the same
timestamped $8\,\mathrm{cm}$ Kimera-VIO pre-mesh and the same VIO
pre-correction and PGMO-optimized trajectories; PGMO uses the official
\texttt{mesh\_trajectory\_deformer} tool~\cite{kimera_pgmo}, whereas
WarpVox applies Stages~1 -- 3.

\begin{table}[t]
\centering
\caption{\textbf{PGMO corrected-surface comparison on ten bags.}}
\label{tab:pgmo_full10}
\scriptsize
\setlength{\tabcolsep}{2.2pt}
\begin{tabular}{lrrrr}
\toprule
Method & F@10 (\%) & F@25 (\%) & F@50 (\%) & Time (s) $\downarrow$ \\
\midrule
Kimera-VIO pre-mesh
& 4.680 & 14.094 & 26.277 & -- \\K
�[Y\�KT�S�ۙK\�������H	��ˎ
�H	�
���N	����H��\����Y�_�HKH‰��
	���L�	�
ˌ��H	�ˍLMH����\�[B�[��X�[\�B�[��X�_B��]�L�X�\��X�\�H	X]�^��_I��۝[�[Y\��\��][ۋ��\����[XZ[���][�	��	�و�[Y\�KT�Sˈ][\�ݙ\���Kѐ
L�B������LI��[H�YX�[���ܜ�X�Y\�\��X�H�ܙH[YH�H	���[Y\����\����X��\]Y[�H\Y\��Y�_���X�ۜ��X�[�	X]�^��_I��]]�H��Y[��	K����X]�^��I
	
��L�X]�^��I�܂��Y�\ߌHKH
H[�ZY[[����Kѐ
L��ܙ\�و	���
�	K�ˌN	I�B�X�][ۈ�\ܝ�H]Z[Y�Y�_��KH
��\\�\�ۋ����X��X�[۞�X�][ۈ�Y_B�X�_��Y��X��X�][۟H\��]\�ۙH\�Yۈ��X�H]H[YHۈۙB���X���\�H�\�H���HXX�وH[��Y��]��WK����Y�[��X�_V�B��[�\�[�—�\[۞�^���X�][ۈو�Y�\ߌHKH
ۈ[��[Y\�KS][H�\�\˟_B�X�[�X��X�][۟B��ܚ\�^�B��][���X����\^̜B��\�^�X�����[[��Y^�_^�B��Y�[��X�[\�^��[�������_B���[B��\�X[�	��L
	JH	���H
	JH	��
L
	JH	��HKH�
�H	���
�H	��
�H�ZY�[B��[[Y\�T�T�	�L���	�����H	�

K�L	�H	�ˌ̌�	�KH���'W7Bc�T�&����b�#srb#b�srbCR�C3bb�b�C�b������FFƖ�W76P��V&W7BԴb�6RG&�7�'@�b��s"b#R�s�bCB�SCbb�CBb�b����W6����ǒE4D`�b�㓃"b#b�CbCR�c�b�b�b������6�W&6R�7W�'FVBE4Db�f���b�SRb#b�cbbCR�3�"b�C�b�C�b�c����&�GF��'V�P��V�G�F'V�'�Х�&&����6��V��v�GF�׵�g76W�'G��67&�G6��RF��W2&RƗ7FVB��ǒf�"F�P�7FvRV�FW"6��&�6���F�RWW"&��6�f&�W232�7F�2&Vf�&R7FvW�B��@��266�&VB��F�R$7W&f6S�F�R��vW"f&�W23��"�B3Bv��7Bf�����B�266�&VB���W6�W2&R�W�G&7FVBg&��F�R�WGWBE4Db�Х�V�G�F&�WР��V&W7BԴb��fW2V6�6��G&��'��G26��6W7B�W�g&�Rw2�6R6�'&V7F�����'6W'fF����FW&�fVBF&vWG2��&�fR��F�&VRb�66�&W2�fW"F��2&6VƖ�R����6�VF��r�c3rBd#R�F�R&�'W7Bc�T�&���&V���2v�F����3"��bgV����W6�$Bd#Rv���R&VGV6��r7FvW�2F��Rg&��2�3#2F��C���2�7FvW�B�VW2F�R&R�W�G&7FVB�W6�v�F����#"�bF�R&�'W7@�7FvW�27W&f6R�B��&�fW2d#R'��#b�fW"�W6����ǒ6��fW'6����@��FF�F����6�7B�b�c%��2ࠤ&W���BF�R7FvR&�F����&V��FVw&F��r��ǒF�R66�VB�W�g&�W2��F�P�6�RFV�66W2&V6�W2C#B���C�c"BBd#R�dS��C32�sE��2�B7F����F�W2�fW"��&FW"�b�v�GVFR���vW"F��v'f���v���R&V����p�C"�CB�2�s�B&V��r�B�F�R66�RF�W&Vf�&R7W�'G2F&vWB6��7G'V7F���'WB��B6���WFRf��V�WG&�2&V'V��Bࠥ�����FV�E�FW�F&g�E4Db6��6�7FV�7��ХvR6��&Rv'f��v�F�6�'&V7FVB��6Rf��&���&V��FVw&F�����6�&V@�f��V�w&�B�G&VF��rV����6FVBf��V�22�W&�vV�v�B�vR&W�'@��6�F�fR�vV�v�B7W�'B��R�vV�v�FVB�66&B6�֖�&�G��bf��V�vV�v�G2��&VfW&V�6R�vV�v�FVBg&VR�76R7W�'B&V6��f�"E������&�'���W2B��@�E4Db�6�v�67W&7���6�����7W�'B�WG6�FRF�RC�W2B�W&���WfV�&�B�fW&vVB�fW"FV�C�#RB��6�V6����G2�v'f��6��WfW2�b���R���e�R�����S�R��B�b�#��R�&W7V7F�fVǒ�6��v��r7V'7F�F��w&VV�V�B��7W�'B�67V�V�FVBvV�v�B��BE4Db6�v�7G'V7GW&Rࠥ�&Vv��F&�WշEХ�6V�FW&��p��6F����FW�F&g�vW76���6V�G&RF&vWB�G&�6fW"66R7GVG�����R���W&��V�F�6WVV�6R��Х��&VǷF#�vW76���66WХ�67&�G6��P��6WF�V�wF���F&6��6W׳"�7GХ�&W6��V&����6��V��v�GF�ײײP��&Vv��F'V�'׶�''''Х�F�'V�P��WF��Bbd#R��R�bdS��R�bFWF���6ҒE�F�v�'&�rBbF��R����֖G'V�P�7�D�&R�6�'&V7F���br�#2b3R�#Rbs2�2b�����v'f��b#2�"bC�S�bC"�SBb#2��7�D�6�'&V7FVB�&V'V��Bb#2�#RbC�bC�3rbE�&��"B�����&�GF��'V�P��V�G�F'V�'�Х�V�G�F&�WР������FV�E�FW�F&g�7&�72�&W&W6V�FF���4Du7��6�FW��W&&�##6vW76��7�66R7GVG��Фǖ��rF�R6�R6�'&V7FVB7W&f6RF&vWG2F�F�RvW76��6V�G&W2�b�7�D���6�FW��VWF�##G7�F��&V6��7G'V7F���&�6�W2F�R6�'&V7FVB��6P�&V'V��B��#��2�v�W&V2F�R&V'V��BF�W2&����FVǒ%�Ƃ�F��0�6��v�R�6WVV�6R&W7V�BWf�VFW2F&vWBG&�6fW&&�ƗG���Ǔ�v'f���&V���2F&vWFVBBE4Db���r�VƖ�W2�
\section{Limitations}
\label{sec:limitations}

WarpVox assumes a static scene, valid geometry, and a locally coherent pre-correction TSDF. It cannot recover geometry absent from the source map or repair incorrectly fused topology, and repeated or nearly overlapping surfaces may still produce an internally consistent but incorrect association. The keyframe-cache memory grows with session length. Finite-resolution TSDF reconstruction cannot exactly reproduce the reintegrated map. The topology-processing stage may remove a small number of faces with unresolved deformation-induced distortion. The attainable accuracy remains limited by residual error in the corrected trajectory, and severe failures require reintegration from the original measurements.
\FloatBarrier
\section{Conclusion}
WarpVox demonstrates that an active TSDF can be corrected after loop closure
without replaying the complete measurement history. Across RGB-D and LiDAR
benchmarks, it remains close to full reintegration; on RGB-D sequences it
achieves a $48\times$ mean update-time speedup. The corrected output is a
native Voxblox volume that immediately resumes fusion and can undergo later
corrections. The Gaussian case study further suggests that the
observation-derived targets are not tied to voxel grids. Code and
configurations will be released.


\bibliographystyle{IEEEtran}
\bibliography{main}

\end{document}